\documentclass{amsart}
% For dealing with references we use the comment environment
\usepackage{verbatim}
\newenvironment{reference}{\comment}{\endcomment}
%\newenvironment{reference}{}{}
\newenvironment{slogan}{\comment}{\endcomment}
\newenvironment{history}{\comment}{\endcomment}

% For commutative diagrams we use Xy-pic
\usepackage[all]{xy}

% We use 2cell for 2-commutative diagrams.
\xyoption{2cell}
\UseAllTwocells

% We use multicol for the list of chapters between chapters
\usepackage{multicol}

% This is generally recommended for better output
\usepackage{lmodern}
\usepackage[T1]{fontenc}

% For cross-file-references

% Package for hypertext links:
\usepackage{hyperref}

% For any local file, say "hello.tex" you want to link to please

% Theorem environments.
%
\theoremstyle{plain}
\newtheorem{theorem}[subsection]{Theorem}
\newtheorem{proposition}[subsection]{Proposition}
\newtheorem{lemma}[subsection]{Lemma}

\theoremstyle{definition}
\newtheorem{definition}[subsection]{Definition}
\newtheorem{example}[subsection]{Example}
\newtheorem{exercise}[subsection]{Exercise}
\newtheorem{situation}[subsection]{Situation}

\theoremstyle{remark}
\newtheorem{remark}[subsection]{Remark}
\newtheorem{remarks}[subsection]{Remarks}

\numberwithin{equation}{subsection}

% Macros
%
\def\lim{\mathop{\mathrm{lim}}\nolimits}
\def\colim{\mathop{\mathrm{colim}}\nolimits}
\def\Spec{\mathop{\mathrm{Spec}}}
\def\Hom{\mathop{\mathrm{Hom}}\nolimits}
\def\Ext{\mathop{\mathrm{Ext}}\nolimits}
\def\SheafHom{\mathop{\mathcal{H}\!\mathit{om}}\nolimits}
\def\SheafExt{\mathop{\mathcal{E}\!\mathit{xt}}\nolimits}
\def\Sch{\mathit{Sch}}
\def\Mor{\mathop{\mathrm{Mor}}\nolimits}
\def\Ob{\mathop{\mathrm{Ob}}\nolimits}
\def\Sh{\mathop{\mathit{Sh}}\nolimits}
\def\NL{\mathop{N\!L}\nolimits}
\def\CH{\mathop{\mathrm{CH}}\nolimits}
\def\proetale{{pro\text{-}\acute{e}tale}}
\def\etale{{\acute{e}tale}}
\def\QCoh{\mathit{QCoh}}
\def\Ker{\mathop{\mathrm{Ker}}}
\def\Im{\mathop{\mathrm{Im}}}
\def\Coker{\mathop{\mathrm{Coker}}}
\def\Coim{\mathop{\mathrm{Coim}}}

% Boxtimes
%
\DeclareMathSymbol{\boxtimes}{\mathbin}{AMSa}{"02}

%
% Macros for moduli stacks/spaces
%
\def\QCohstack{\mathcal{QC}\!\mathit{oh}}
\def\Cohstack{\mathcal{C}\!\mathit{oh}}
\def\Spacesstack{\mathcal{S}\!\mathit{paces}}
\def\Quotfunctor{\mathrm{Quot}}
\def\Hilbfunctor{\mathrm{Hilb}}
\def\Curvesstack{\mathcal{C}\!\mathit{urves}}
\def\Polarizedstack{\mathcal{P}\!\mathit{olarized}}
\def\Complexesstack{\mathcal{C}\!\mathit{omplexes}}
% \Pic is the operator that assigns to X its picard group, usage \Pic(X)
% \Picardstack_{X/B} denotes the Picard stack of X over B
% \Picardfunctor_{X/B} denotes the Picard functor of X over B
\def\Pic{\mathop{\mathrm{Pic}}\nolimits}
\def\Picardstack{\mathcal{P}\!\mathit{ic}}
\def\Picardfunctor{\mathrm{Pic}}
\def\Deformationcategory{\mathcal{D}\!\mathit{ef}}

\setlength{\emergencystretch}{3em}
\begin{document}
\title[Stacks Project, Tag 0H7R]{336 --- Descent of projectivity through finite injective ring extensions}
\maketitle
\noindent Copyright (C) 2005--2025 Johan de Jong. Stacks Project authors. Source: \href{https://stacks.math.columbia.edu/tag/0H7R}{Tag 0H7R}.

\noindent Editorial difficulty note (not author text). Difficulty H2: The proof changes to a faithfully flat finite locally free extension where the integral algebra splits into polynomial-root pieces. It patches projectivity across quotient ideals, controls the nilpotent intersection, and combines descended flatness with lifting of projectives; the extension is not assumed flat..

\noindent Editorial provenance note (not author text). History: excerpt from revision \texttt{a0444\allowbreak 6e57e\allowbreak c1fbc\allowbreak 252a8\allowbreak 71afc\allowbreak ec775\allowbreak 2fb28\allowbreak 07b14}, more-algebra.tex, lines 4721--4764. Reference presentation was adapted; original-author context and core support blocks were retained or added. The mathematical wording is unchanged. Licensed under GNU FDL 1.2 or later; no invariant sections or cover texts. See ../licenses/COPYING. Context blocks reproduce author definitions and supporting results; they are not additional counted problems.

\begin{lemma}
\label{lemma-have-one-root}
Let $R$ be a ring. Let $P(T)$ be a monic polynomial with coefficients
in $R$. Let $\alpha \in R$ be such that $P(\alpha) = 0$. Then
$P(T) = (T - \alpha)Q(T)$ for some monic polynomial $Q(T) \in R[T]$.
\end{lemma}

\begin{proof}
By induction on the degree of $P$. If $\deg(P) = 1$, then
$P(T) = T - \alpha$ and the result is true. If $\deg(P) > 1$, then
we can write $P(T) = (T - \alpha)Q(T) + r$ for some polynomial
$Q \in R[T]$ of degree $< \deg(P)$ and some $r \in R$ by long
division. By assumption $0 = P(\alpha) = (\alpha - \alpha)Q(\alpha) + r = r$
and we conclude that $r = 0$ as desired.
\end{proof}


\begin{lemma}
\label{lemma-adjoin-one-root}
Let $R$ be a ring. Let $P(T)$ be a monic polynomial with coefficients
in $R$. There exists a finite free ring map $R \to R'$ such that
$P(T) = (T - \alpha)Q(T)$ for some $\alpha \in R'$ and some
monic polynomial $Q(T) \in R'[T]$.
\end{lemma}

\begin{proof}
Write $P(T) = T^d + a_1T^{d - 1} + \ldots + a_0$.
Set $R' = R[x]/(x^d + a_1x^{d - 1} + \ldots + a_0)$.
Set $\alpha$ equal to the congruence class of $x$.
Then it is clear that $P(\alpha) = 0$. Thus we win by
Lemma \ref{lemma-have-one-root}.
\end{proof}


\begin{lemma}
\label{lemma-finite-split}
Let $R \to S$ be a finite ring map.
There exists a finite free ring extension $R \subset R'$ such
that $S \otimes_R R'$ is a quotient of a ring of the form
$$
R'[T_1, \ldots, T_n]/(P_1(T_1), \ldots, P_n(T_n))
$$
with $P_i(T) = \prod_{j = 1, \ldots, d_i} (T - \alpha_{ij})$ for some
$\alpha_{ij} \in R'$.
\end{lemma}

\begin{proof}
Let $x_1, \ldots, x_n \in S$ be generators of $S$ over $R$.
For each $i$ we can choose a monic polynomial $P_i(T) \in R[T]$
such that $P_i(x_i) = 0$ in $S$, see
Algebra, Lemma \href{https://stacks.math.columbia.edu/tag/00GK}{lemma finite is integral (Tag 00GK)}.
Say $\deg(P_i) = d_i$. By
Lemma \ref{lemma-adjoin-one-root}
(applied $\sum d_i$ times) there exists a finite free ring
extension $R \subset R'$ such that each $P_i$ splits completely:
$$
P_i(T) = \prod\nolimits_{j = 1, \ldots, d_i} (T - \alpha_{ij})
$$
for certain $\alpha_{ik} \in R'$. Let
$R'[T_1, \ldots, T_n] \to S \otimes_R R'$ be the $R'$-algebra map
which maps $T_i$ to $x_i \otimes 1$. As this maps $P_i(T_i)$ to zero,
this induces the desired surjection.
\end{proof}


\begin{lemma}
\label{lemma-split-image}
Let $R$ be a ring.
Let $S = R[T_1, \ldots, T_n]/J$.
Assume $J$ contains elements of the form $P_i(T_i)$
with $P_i(T) = \prod_{j = 1, \ldots, d_i} (T - \alpha_{ij})$ for some
$\alpha_{ij} \in R$. For $\underline{k} = (k_1, \ldots, k_n)$
with $1 \leq k_i \leq d_i$ consider the ring map
$$
\Phi_{\underline{k}} : R[T_1, \ldots, T_n] \to R,
\quad
T_i \longmapsto \alpha_{ik_i}
$$
Set $J_{\underline{k}} = \Phi_{\underline{k}}(J)$.
Then the image of $\Spec(S) \to \Spec(R)$ is equal to
$V(\bigcap J_{\underline{k}})$.
\end{lemma}

\begin{proof}
This lemma proves itself. Hint:
$V(\bigcap J_{\underline{k}}) = \bigcup V(J_{\underline{k}})$.
\end{proof}


\begin{lemma}
\label{lemma-descent-flatness-injective-finite-Noetherian-rings}
Let $R \to S$ be a finite injective homomorphism of Noetherian rings.
Let $M$ be an $R$-module. If $M \otimes_R S$ is a flat $S$-module,
then $M$ is a flat $R$-module.
\end{lemma}

\begin{proof}
Let $M$ be an $R$-module such that $M \otimes_R S$ is flat over $S$. By
Algebra, Lemma \href{https://stacks.math.columbia.edu/tag/00HJ}{lemma flatness descends (Tag 00HJ)}
in order to prove that $M$ is flat we may replace $R$ by any faithfully flat
ring extension. By
Lemma \ref{lemma-finite-split}
we can find a finite locally free ring extension $R \subset R'$ such
that $S' = S \otimes_R R' = R'[T_1, \ldots, T_n]/J$ for some ideal
$J \subset R'[T_1, \ldots, T_n]$ which contains the  elements of the form
$P_i(T_i)$ with $P_i(T) = \prod_{j = 1, \ldots, d_i} (T - \alpha_{ij})$
for some $\alpha_{ij} \in R'$. Note that $R'$ is Noetherian
and that $R' \subset S'$ is a finite extension of rings. Hence we may
replace $R$ by $R'$ and assume that $S$ has a presentation as in
Lemma \ref{lemma-split-image}.
Note that $\Spec(S) \to \Spec(R)$ is surjective, see
Algebra, Lemma \href{https://stacks.math.columbia.edu/tag/00GQ}{lemma integral overring surjective (Tag 00GQ)}.
Thus, using
Lemma \ref{lemma-split-image}
we conclude that $I = \bigcap J_{\underline{k}}$ is an ideal
such that $V(I) = \Spec(R)$. This means that
$I \subset \sqrt{(0)}$, and since $R$ is Noetherian that $I$
is nilpotent. The maps $\Phi_{\underline{k}}$ induce commutative
diagrams
$$
\xymatrix{
S \ar[rr] & & R/J_{\underline{k}} \\
& R \ar[lu] \ar[ru]
}
$$
from which we conclude that $M/J_{\underline{k}}M$ is flat over
$R/J_{\underline{k}}$. By
Lemma \ref{lemma-intersection-flat}
we see that $M/IM$ is flat over $R/I$. Finally, applying
Algebra, Lemma \ref{algebra-lemma-lift-flatness}
we conclude that $M$ is flat over $R$.
\end{proof}


\begin{lemma}
\label{lemma-intersection-flat}
Let $R$ be a ring. Let $M$ be an $R$-module. Let $I_1$, $I_2$ be ideals of $R$.
If $M/I_1M$ is flat over $R/I_1$ and $M/I_2M$ is flat over $R/I_2$,
then $M/(I_1 \cap I_2)M$ is flat over $R/(I_1 \cap I_2)$.
\end{lemma}

\begin{proof}
By replacing $R$ with $R/(I_1 \cap I_2)$ and $M$ by $M/(I_1 \cap I_2)M$
we may assume that $I_1 \cap I_2 = 0$. Let $J \subset R$ be an ideal.
To prove that $M$ is flat over $R$ we have to show that
$J \otimes_R M \to M$ is injective, see
Algebra, Lemma \href{https://stacks.math.columbia.edu/tag/00HD}{lemma flat (Tag 00HD)}.
By flatness of $M/I_1M$ over $R/I_1$ the map
$$
J/(J \cap I_1) \otimes_R M =
(J + I_1)/I_1 \otimes_{R/I_1} M/I_1M
\longrightarrow M/I_1M
$$
is injective. As $0 \to (J \cap I_1) \to J \to J/(J \cap I_1) \to 0$
is exact we obtain a diagram
$$
\xymatrix{
(J \cap I_1) \otimes_R M \ar[r] \ar[d] &
J \otimes_R M \ar[r] \ar[d] &
J/(J \cap I_1) \otimes_R M \ar[r] \ar[d] & 0 \\
M \ar@{=}[r] &
M \ar[r] &
M/I_1M
}
$$
hence it suffices to show that $(J \cap I_1) \otimes_R M \to M$ is
injective. Since $I_1 \cap I_2 = 0$ the ideal $J \cap I_1$
maps isomorphically to an ideal $J' \subset R/I_2$ and we see that
$(J \cap I_1) \otimes_R M = J' \otimes_{R/I_2} M/I_2M$. By flatness
of $M/I_2M$ over $R/I_2$ the map $J' \otimes_{R/I_2} M/I_2M \to M/I_2M$
is injective, which clearly implies that $(J \cap I_1) \otimes_R M \to M$ is
injective.
\end{proof}


\begin{lemma}
\label{algebra-lemma-lift-flatness}
Let $\varphi : R \to R'$ be a ring map.
Let $I \subset R$ be an ideal.
Let $M$ be an $R$-module.
Assume
\begin{enumerate}
\item $I$ is nilpotent,
\item $R \to R'$ is injective,
\item $M/IM$ is flat over $R/I$, and
\item $R' \otimes_R M$ is flat over $R'$.
\end{enumerate}
Then $M$ is flat over $R$.
\end{lemma}

\begin{proof}
Define inductively $I_1 = I$ and $I_{n + 1} = \varphi^{-1}(\varphi(I_n)^2R')$
for $n \geq 1$. Note that by
Lemma \ref{algebra-lemma-prepare-lift-flatness}
we find that $M/I_nM$ is flat over $R/I_n$ for each $n \geq 1$.
It is clear that $\varphi(I_{n + 1}) \subset \varphi(I)^{2^n}R'$. Since
$I$ is nilpotent we see that $\varphi(I_n) = 0$ for some $n$. As
$\varphi$ is injective we conclude that $I_n = 0$ for some $n$ and
we win.
\end{proof}


\begin{lemma}
\label{algebra-lemma-prepare-lift-flatness}
Let $\varphi : R \to R'$ be a ring map.
Let $I \subset R$ be an ideal.
Let $M$ be an $R$-module.
Assume
\begin{enumerate}
\item $M/IM$ is flat over $R/I$, and
\item $R' \otimes_R M$ is flat over $R'$.
\end{enumerate}
Set $I_2 = \varphi^{-1}(\varphi(I^2)R')$.
Then $M/I_2M$ is flat over $R/I_2$.
\end{lemma}

\begin{proof}
We may replace $R$, $M$, and $R'$ by $R/I_2$, $M/I_2M$, and
$R'/\varphi(I)^2R'$. Then $I^2 = 0$ and $\varphi$ is injective. By
Lemma \href{https://stacks.math.columbia.edu/tag/051C}{lemma what does it mean (Tag 051C)}
and the fact that $I^2 = 0$ it suffices to prove that
$\text{Tor}^R_1(R/I, M) = K = \Ker(I \otimes_R M \to M)$ is zero.
Set $M' = M \otimes_R R'$ and $I' = IR'$.
By assumption the map $I' \otimes_{R'} M' \to M'$ is injective.
Hence $K$ maps to zero in
$$
I' \otimes_{R'} M' = I' \otimes_R M = I' \otimes_{R/I} M/IM.
$$
Then $I \to I'$ is an injective map of $R/I$-modules.
Since $M/IM$ is flat over $R/I$ the map
$$
I \otimes_{R/I} M/IM \longrightarrow I' \otimes_{R/I} M/IM
$$
is injective. This implies that $K$ is zero in
$I \otimes_R M = I \otimes_{R/I} M/IM$ as desired.
\end{proof}


\begin{lemma}
\label{algebra-lemma-projective-modulo-two-ideals}
Let $R$ be a ring. Let $I, J \subset R$ be ideals such that $I \cap J = 0$.
Let $P$ be an $R$-module such that $P/IP$ is a projective $R/I$-module and
$P/JP$ is a projective $R/J$-module. Then $P$ is a projective $R$-module.
\end{lemma}

\begin{proof}
Choose a surjection $p : F \to P$ where $F$ is a free $R$-module.
Since $P/IP$ is a projective $R/I$-module, we can choose a map
$f : P/IP \to F/IF$ which is a right inverse to $p \bmod I$.
Consider the map $q : F/JF \to F/(I + J)F \times_{P/(I + J)P} P/JP$
induced by the quotient map $F/JF \to F/(I + J)F$ and $p$.
Note that $q$ is a surjective map of $R/J$-modules (small detail omitted).
Consider the map $f' : P/JP \to F/(I + J)F \times_{P/(I + J)P} P/JP$
induced by $f$ and the identity map.
Since $P/JP$ is a projective $R/J$-module and $q$ surjective,
we can choose a map $g : P/JP \to F/JF$ such that $q \circ g = f'$.
Then $f \bmod I + J = g \bmod I + J$. Since $F = F/JF \times_{F/(I + J)F} F/IF$
because $I \cap J = 0$, we conclude that we obtain a well defined homomorphism
$h = (f, g) : P \to F$ of $R$-modules. The map $a = p \circ h : P \to P$
reduces to the identity modulo $I$ and modulo $J$. Then $a$ also induces
the identity map on the submodule $IP$: if $x = \sum i_\alpha x_\alpha$
with $i_\alpha \in I$ and $x_\alpha \in P$, then
$a(x) = \sum i_\alpha a(x_\alpha) = \sum i_\alpha x_\alpha = x$
because $a(x_\alpha) = x_\alpha + j_\alpha$ with $j_\alpha \in J$
and $i_\alpha j_\alpha = 0$. Thus $a : P \to P$ acts as the identity
on $IP$ and on $P/IP$ and we conclude that $a$ is an automorphism of $P$
(small detail omitted). Thus $P$ is a summand of $F$, whence projective.
\end{proof}


\begin{lemma}
\label{algebra-lemma-lift-projective}
Let $R$ be a ring.
Let $I \subset R$ be an ideal.
Let $M$ be an $R$-module.
Assume
\begin{enumerate}
\item $I$ is nilpotent,
\item $M/IM$ is a projective $R/I$-module,
\item $M$ is a flat $R$-module.
\end{enumerate}
Then $M$ is a projective $R$-module.
\end{lemma}

\begin{proof}
By Lemma \ref{algebra-lemma-lift-projective-module} we can find a projective
$R$-module $P$ and an isomorphism $P/IP \to M/IM$. We are going to show
that $M$ is isomorphic to $P$ which will finish the proof. Because $P$
is projective we can lift the map $P \to P/IP \to M/IM$ to an $R$-module
map $P \to M$ which is an isomorphism modulo $I$. Since $I^n = 0$
for some $n$, we can use the filtrations
\begin{align*}
0 = I^nM \subset I^{n - 1}M \subset \ldots \subset IM \subset M \\
0 = I^nP \subset I^{n - 1}P \subset \ldots \subset IP \subset P
\end{align*}
to see that it suffices to show that the induced maps
$I^aP/I^{a + 1}P \to I^aM/I^{a + 1}M$ are bijective. Since both $P$
and $M$ are flat $R$-modules we can identify this with the map
$$
I^a/I^{a + 1} \otimes_{R/I} P/IP
\longrightarrow
I^a/I^{a + 1} \otimes_{R/I} M/IM
$$
induced by $P \to M$. Since we chose $P \to M$ such that the induced
map $P/IP \to M/IM$ is an isomorphism, we win.
\end{proof}


\begin{lemma}
\label{algebra-lemma-lift-projective-module}
Let $R$ be a ring. Let $I \subset R$ be a nilpotent ideal. Let
$\overline{P}$ be a projective $R/I$-module. Then there exists a
projective $R$-module $P$ such that $P/IP \cong \overline{P}$.
\end{lemma}

\begin{proof}
By Lemma \href{https://stacks.math.columbia.edu/tag/05CF}{lemma characterize projective (Tag 05CF)}
we can choose a set $A$ and a direct sum decomposition
$\bigoplus_{\alpha \in A} R/I = \overline{P} \oplus \overline{K}$
for some $R/I$-module $\overline{K}$. Write $F = \bigoplus_{\alpha \in A} R$
for the free $R$-module on $A$. Choose a lift
$p : F \to F$ of the projector $\overline{p}$ associated to
the direct summand $\overline{P}$ of $\bigoplus_{\alpha \in A} R/I$.
Note that $p^2 - p \in \text{End}_R(F)$ is a nilpotent
endomorphism of $F$ (as $I$ is nilpotent and the matrix entries of
$p^2 - p$ are in $I$; more precisely, if $I^n = 0$, then $(p^2 - p)^n = 0$).
Hence by Lemma \href{https://stacks.math.columbia.edu/tag/05BU}{lemma lift idempotents noncommutative (Tag 05BU)}
we can modify our choice of $p$ and assume that $p$ is a projector.
Set $P = \Im(p)$.
\end{proof}


\begin{lemma}
\label{lemma-descent-projective-injective-finite-Noetherian-rings}
Let $R \to S$ be a finite injective homomorphism of Noetherian rings.
Let $P$ be an $R$-module. If $P \otimes_R S$ is a projective $S$-module,
then $P$ is a projective $R$-module.
\end{lemma}

\begin{proof}
Let $P$ be an $R$-module such that $P \otimes_R S$ is projective over $S$. By
Algebra, Theorem \href{https://stacks.math.columbia.edu/tag/05A9}{theorem ffdescent projectivity (Tag 05A9)}
in order to prove that $P$ is projective we may replace $R$
by any faithfully flat ring extension. By
Lemma \ref{lemma-finite-split}
we can find a finite locally free ring extension $R \subset R'$ such
that $S' = S \otimes_R R' = R'[T_1, \ldots, T_n]/J$ for some ideal
$J \subset R'[T_1, \ldots, T_n]$ which contains the  elements of the form
$P_i(T_i)$ with $P_i(T) = \prod_{j = 1, \ldots, d_i} (T - \alpha_{ij})$
for some $\alpha_{ij} \in R'$. Note that $R'$ is Noetherian
and that $R' \subset S'$ is a finite extension of rings. Hence we may
replace $R$ by $R'$ and assume that $S$ has a presentation as in
Lemma \ref{lemma-split-image}.
Note that $\Spec(S) \to \Spec(R)$ is surjective, see
Algebra, Lemma \href{https://stacks.math.columbia.edu/tag/00GQ}{lemma integral overring surjective (Tag 00GQ)}.
Thus, using Lemma \ref{lemma-split-image}
we conclude that $I = \bigcap J_{\underline{k}}$ is an ideal
such that $V(I) = \Spec(R)$. This means that
$I \subset \sqrt{(0)}$, and since $R$ is Noetherian that $I$
is nilpotent. The maps $\Phi_{\underline{k}}$ induce commutative
diagrams
$$
\xymatrix{
S \ar[rr] & & R/J_{\underline{k}} \\
& R \ar[lu] \ar[ru]
}
$$
from which we conclude that $P/J_{\underline{k}}P$ is projective over
$R/J_{\underline{k}}$. By
Algebra, Lemma \ref{algebra-lemma-projective-modulo-two-ideals}
we see that $P/IP$ is projective over $R/I$. Since $P$ is flat
over $R$ by
Lemma \ref{lemma-descent-flatness-injective-finite-Noetherian-rings},
we may apply Algebra, Lemma \ref{algebra-lemma-lift-projective} to
conclude that $P$ is projective over $R$.
\end{proof}
\end{document}
